\documentclass[11pt,a4paper]{article}
\usepackage[T1]{fontenc}
\usepackage{amsmath,amssymb,amsthm,mathtools}
\usepackage[margin=27mm]{geometry}
\usepackage{lmodern,microtype,booktabs}
\usepackage[hidelinks]{hyperref}
\usepackage{xurl}
\hypersetup{pdftitle={Weighted Rank Five and Covers with a Two Vertex Shore},pdfauthor={Research note prepared for Vladimir Reshetnikov with OpenAI}}
\setlength{\parskip}{2pt}
\newtheorem{theorem}{Theorem}[section]
\newtheorem{lemma}[theorem]{Lemma}
\newtheorem{corollary}[theorem]{Corollary}
\theoremstyle{remark}\newtheorem{remark}[theorem]{Remark}
\newcommand{\ULC}{\operatorname{ULC}}
\newcommand{\HH}{\mathbb H}
\newcommand{\e}{\mathrm e}
\title{Weighted Rank Five\newline and Covers with a Two Vertex Shore}
\author{Research note prepared for Vladimir Reshetnikov with OpenAI}
\date{1 October 2026}
\begin{document}
\maketitle
\begin{abstract}
Every bipartite matching-support polynomial with a minimum vertex cover meeting one shore in at most two vertices is ultra-log-concave after normalization by its matching rank, for arbitrary positive activities on both shores. The main construction replaces a displayed $2+q$ cover by a rank-$(q+2)$ represented matroid whose Lorentzian basis polynomial specializes exactly to the homogenized support polynomial. Consequently weighted rank normalization holds through rank five, and known rank-six Hall obstructions show that the first weighted failure rank is exactly six. We also prove a sharp all-core boundary inequality with constant three, give independent Lorentzian transfer proofs for star cores, and derive quantitative eventual strictness under two natural activity scalings. Full proofs, finite exact coefficient certificates, sum-of-squares identities and independent verification scripts are included.
\end{abstract}

\paragraph{Scope and status.}
The rank-five weighted question posed in the inspected ProveIt draft is resolved here. At higher rank the theorem covers every minimum cover with a shore of size at most two; the unrestricted unit-weight and one-shore-weighted problems at higher rank are not addressed here. These are proofs in ordinary mathematics, not Lean-certified or independently refereed results. The work extends the inspected rank manuscript and the separate rank-four and full-Hall classification reports \cite{Rank,Four,Hall}; global priority is not claimed.

\section{Definitions and main conclusions}
For a finite bipartite graph $G=(X,Y;E)$ and positive vertex activities $u_x,v_y$, write
\[
 p_G(t)=\sum_{k=0}^{r}a_kt^k,
 \qquad
 a_k=\sum_{\substack{I\subseteq X,\ J\subseteq Y\\|I|=|J|=k\\G[I,J]\text{ perfectly matchable}}}u_Iv_J,
 \qquad r=\nu(G).
\]
Each endpoint pair is counted once. In particular $p_{K_{2,2}}(t)=1+4t+t^2$ at unit activities. The polynomial is $\ULC_D$, for $D\ge\deg p_G$, if its zero-padded sequence $a_k/\binom Dk$ is log-concave and has no internal zeros. Equivalently, for these support polynomials,
\begin{equation}\label{eq:gap}
 \Delta_k^{(D)}:=k(D-k)a_k^2-(k+1)(D-k+1)a_{k-1}a_{k+1}\ge0.
\end{equation}
Strict $\ULC_D$ means all $D-1$ displayed gaps are positive.

A displayed $2+q$ cover is $A\sqcup B$, where $A=\{A_1,A_2\}\subseteq X$ and $|B|=q$, $B\subseteq Y$. Put $L=X\setminus A$ and $R=Y\setminus B$. There are no $L$--$R$ edges. The induced graph $G[A,B]$ is called the core. We call the displayed cover \emph{genuine} when $\nu(G)=q+2$.

\begin{theorem}[A two-vertex cover shore]\label{thm:main}
If $G$ has a displayed $2+q$ cover, fix positive activities $v$ on its right shore. The homogeneous left marginal
\[
 \sum_{\substack{I\subseteq X\\|I|\le q+2}}\left(\sum_{\substack{J\subseteq Y\\(I,J)\text{ feasible}}}v_J\right)z_I s^{q+2-|I|}
\]
is Lorentzian. In particular, the degree-$(q+2)$ homogenization of its weighted matching-support polynomial is Lorentzian. Thus $p_G$ is $\ULC_{q+2}$ for every positive activity assignment. If the displayed cover is minimum, the normalization is by the actual matching rank.
\end{theorem}

\begin{corollary}[Weighted rank five and the first failure rank]\label{cor:rankfive}
Every bipartite matching-support polynomial of matching rank $r\le5$ is $\ULC_r$ for arbitrary positive activities on both shores. The smallest rank admitting a weighted counterexample is exactly six. At every rank $r\ge6$ there is a weighted counterexample.
\end{corollary}
The positive conclusion follows from Theorem~\ref{thm:main}, the pure and singleton-side cover cases proved below, and K\H{o}nig's theorem. The negative construction is given in Section~\ref{sec:negative}. For nonnegative activities, delete zero-activity vertices first; the same conclusion holds with $r$ interpreted as the surviving degree.

\begin{theorem}[Forced-sector and strict-scaling refinements]\label{thm:forcedleft}
For any displayed $2+q$ cover with $q\ge1$, the coefficient of the product of the two $A$ activities in $p_G$ is $t^2C(t)$, where $C$ is $\ULC_q$. For a genuine cover, common scaling of only the two $A$ activities makes $p_G$ strictly $\ULC_{q+2}$ for all sufficiently large scale factors.

For a rank-five $2+3$ cover whose core has matching rank two, common scaling of all five cover activities also makes all four gaps strictly positive eventually. An explicit coefficient bound gives a sufficient threshold for both scalings.
\end{theorem}
Every fixed minimum-cover scaling at rank at most five is therefore $\ULC$ at \emph{every} positive scale factor. At each rank at least six, some fixed graph and minimum cover fail for every sufficiently large common cover scale, so the universal scaling threshold is sharp as well.

\section{Lorentzian input and singleton covers}
We use the following published facts from Br\"and\'en--Huh \cite{BH}: matroid basis polynomials are Lorentzian (Theorem 3.10); partial derivatives, nonnegative linear substitutions, products and coefficientwise limits preserve this class (Corollary 2.11, Theorem 2.10, Corollary 2.32, and closedness); homogeneous stability preservers that preserve nonnegative coefficients preserve Lorentzian polynomials (Theorem 3.4); and multiaffine truncation and polarization preserve the class (Corollary 3.5 and Proposition 3.1). A homogeneous bivariate polynomial of degree $D$ is Lorentzian exactly when its coefficients form a $\ULC_D$ sequence (Example 2.26). The zero polynomial is included by closure.

\begin{lemma}[Weighted smaller-shore inequality]\label{lem:small}
Every bipartite support polynomial is $\ULC_{\min(|X|,|Y|)}$.
\end{lemma}
\begin{proof}
Add a private element $\bar x$ for each $x\in X$. The transversal matroid on $\bar X\sqcup Y$, matched into $X$, has bases exactly $\overline{X\setminus I}\sqcup J$ for feasible $(I,J)$. If $F$ is its basis polynomial, then
\[
 \left(\prod_{x\in X}u_x\right)
 F\big((s/u_x)_{x\in X},(tv_y)_{y\in Y}\big)
 =\sum_k a_k s^{|X|-k}t^k.
\]
The Lorentzian facts give $\ULC_{|X|}$; transpose the graph for the other shore. This is the standard bimatroid specialization \cite{RU}.
\end{proof}

\begin{lemma}[Singleton-side minimum covers]\label{lem:singleton}
If a rank-$r$ graph has a minimum cover entirely on one shore or split as $1+(r-1)$, its support polynomial is $\ULC_r$ for all positive activities.
\end{lemma}
\begin{proof}
For a pure cover, delete isolated vertices on that shore and apply Lemma~\ref{lem:small}. For a cover $\{x_0\}\sqcup B$, every nonisolated right vertex outside $B$ is a leaf at $x_0$. A feasible support contains at most one such leaf; replacing all of them by one leaf with activity equal to their sum preserves the entire support polynomial. There is at least one such leaf, since the rank exceeds $|B|$. The compressed graph has exactly $r$ active right vertices and the same polynomial, so Lemma~\ref{lem:small} applies. This is the exact leaf-compression argument of \cite{Leaf}.
\end{proof}

\section{The exact two-row formula}
Let the two $A$ activities be $x,y$, and let $H=G[L,B]$. For a virtual left vertex with neighborhood $S\subseteq B$ and activity one, define $g_S$ by
\[
 p_H=f,\qquad p_{H+S}=f+t g_S.
\]
If $S_1,S_2$ are the actual core-row neighborhoods, define $h$ by
\[
 p_{H+A}=f+t(xg_1+yg_2)+xy t^2h,
 \qquad g_i=g_{S_i},\quad g_U=g_{S_1\cup S_2}.
\]
The feasible support set for $g_U$ is the union of those for $g_1,g_2$: in a matching the virtual row uses one edge, and that edge belongs to at least one of the original neighborhoods.

For $R$, let $a,b,c$ be the sums of activities in classes with neighborhoods $\{A_1\}$, $\{A_2\}$ and $A$, respectively. If the individual common-class activities are $r_1,\ldots,r_m$, put
\[
 d=\sum_{i<j}r_i r_j,\qquad Q=ab+ac+bc+d,
 \qquad 0\le d\le c^2/2.
\]
Then the exact full polynomial is
\begin{align}\label{eq:union}
 p_G={}&f\big[1+(x(a+c)+y(b+c))t+xyQt^2\big]
       +t(xg_1+yg_2)\notag\\
 &+xy t^2\big(h+a g_2+b g_1+c g_U\big).
\end{align}
To prove this, condition on the selected $A$ vertices and then on the number of selected $R$ vertices. With both $A$ selected, two selected $R$ contribute $Qf$; one singleton-class vertex forces the corresponding $A$ and gives $ag_2$ or $bg_1$; one common-class vertex leaves the union-neighborhood alternative and gives $cg_U$; no selected $R$ gives $h$. These cases have disjoint endpoint sets. In particular the common term is $cg_U$, not $c(g_1+g_2)$.

\section{The full rank compression matroid lift}\label{sec:fulllift}
We prove Theorem~\ref{thm:main}. First suppose $Q>0$. The same construction includes $q=0$, with no original $B$ slots. Form a transversal matroid matched into the $q+2$ slots
\[
 B\sqcup\{e_1,e_2\}.
\]
Its elements and allowed slots are as follows:
\begin{itemize}
\item a private element $\bar b$ for each $b\in B$, allowed only in slot $b$;
\item an element for each $l\in L$, allowed in its original neighborhood in $B$;
\item two augmented core elements $\widehat A_i$, allowed in $S_i\cup\{e_i\}$;
\item two plane elements $P_1,P_2$, allowed only in $e_1,e_2$, respectively;
\item one common element $Z_j$ for each common right-exterior vertex, allowed in both $e_1,e_2$.
\end{itemize}
The rank is $q+2$, since all private $B$ elements together with $P_1,P_2$ form a basis. Every basis is counted once, regardless of how many injections into the slots it admits. We call $P_1,P_2,Z_j$ the plane elements, as they use only the two additional slots.

Give the private $B$ variables the values $s/v_b$, and the exterior-left variables $tu_l$. Keep core variables $X,Y$, and specialize the plane variables to
\begin{equation}\label{eq:planeweights}
 P_1:\frac{bs}{Q},\qquad P_2:\frac{as}{Q},\qquad
 Z_j:\frac{r_js}{Q}.
\end{equation}
Multiply this basis specialization by $Q\prod_{b\in B}v_b$. It is Lorentzian by the matroid and nonnegative-substitution theorems.

We now compute it exactly. Put
\[
 f_h=s^qf(t/s),\quad G_i=s^{q-1}g_i(t/s),\quad
 G_U=s^{q-1}g_U(t/s),\quad H_h=s^{q-2}h(t/s),
\]
where $G_i=G_U=0$ is defined directly for $q=0$, and $H_h=0$ is defined directly for $q<2$; the displayed homogenizing formulas are used only for nonnegative exponents. Classify a basis by the number $k$ of core elements and $\ell$ of plane elements. Other elements use only the $q$ original slots, so $k+\ell\ge2$ and $0\le k,\ell\le2$. There are six cases.
\begin{enumerate}
\item $(k,\ell)=(0,2)$. Any pair of distinct plane elements can fill the two extra slots, and the other $q$ elements fill the $B$ slots. The sum of plane-pair activities is
\[
 \frac{s^2}{Q^2}\left(ab+(a+b)c+\sum_{i<j}r_ir_j\right)=\frac{s^2}{Q}.
\]
After the exterior factor $Q$, the contribution is $s^2f_h$.
\item $(k,\ell)=(1,1)$. With core $\widehat A_1$, the plane element must be able to occupy $e_2$, hence is $P_2$ or a $Z_j$. Its total activity is $(a+c)s/Q$. The remaining $q$ elements fill the $B$ slots, giving $Xs(a+c)f_h$. With $\widehat A_2$, the contribution is $Ys(b+c)f_h$.
\item $(k,\ell)=(1,2)$. The plane pair fills both extra slots, forcing the core element into $B$. The remaining support uses row $S_i$ and $q-1$ ordinary elements, with profile $G_i$. The contribution is $Xs^2G_1+Ys^2G_2$.
\item $(k,\ell)=(2,0)$. The two core elements occupy their dedicated extra slots. The other $q$ elements fill $B$, contributing $QXYf_h$.
\item $(k,\ell)=(2,1)$. The element $P_1$ occupies $e_1$, forcing $\widehat A_2$ into $e_2$ and $\widehat A_1$ into $B$; hence its profile is $G_1$ and its activity coefficient is $b$. Similarly $P_2$ gives $G_2$ with coefficient $a$. A common $Z_j$ can occupy either extra slot, so feasibility holds exactly when either core row completes the remaining $B$-slot support. This is the union profile $G_U$, counted once even if both alternatives work. Thus the contribution is $XYs(bG_1+aG_2+cG_U)$.
\item $(k,\ell)=(2,2)$. The two plane elements fill the extra slots, forcing both core elements into $B$ together with $q-2$ other elements. The contribution is $XYs^2H_h$; this case is absent for $q<2$.
\end{enumerate}
Consequently the constructed Lorentzian polynomial is
\begin{align}\label{eq:fulllift}
 \mathcal L(s,t,X,Y)={}&s^2f_h+sX\big[(a+c)f_h+sG_1\big]
 +sY\big[(b+c)f_h+sG_2\big]\notag\\
 &+XY\big[Qf_h+s(bG_1+aG_2+cG_U)+s^2H_h\big].
\end{align}
Set $X=xt$ and $Y=yt$. Equation~\eqref{eq:union} identifies \eqref{eq:fulllift} with $s^{q+2}p_G(t/s)$, proving the bivariate assertion for $Q>0$. More generally, keep each exterior-left element variable as its own $z_l$ rather than $tu_l$, and set $X=z_{A_1}$, $Y=z_{A_2}$. The same six basis cases give exactly the homogeneous left marginal in Theorem~\ref{thm:main}. Every plane and private variable remains a nonnegative multiple of $s$, so this full marginal is Lorentzian. No analogous opposite-marginal assertion is inferred when $q>2$.

The six cases are disjoint classes of matroid bases. In particular the union case does not count injections or matching allocations. The construction also has a generic matrix representation obtained by giving each allowed slot incidence its own indeterminate. Its six block-determinant expansions provide a second interpretation of the same identity, including dependent or zero original core-row vectors.

If $Q=0$, add a singleton right vertex at each $A_i$ with activity $\varepsilon>0$. The perturbed graph still has a displayed $2+q$ cover, its $Q$ is positive, and its coefficients converge to those of the original graph as $\varepsilon\downarrow0$. Closedness gives the result. Other zero activities are handled by continuity at fixed order $q+2$; private reciprocals are used only before taking the positive-activity limit.

\begin{corollary}\label{cor:coverobstruction}
Any failure of weighted rank-normalized ultra-log-concavity must have every minimum vertex cover meet both shores in at least three vertices.
\end{corollary}
\begin{proof}
A pure or singleton-side minimum cover is covered by Lemma~\ref{lem:singleton}; a cover meeting one shore in exactly two vertices is covered by Theorem~\ref{thm:main}.
\end{proof}

\begin{remark}
The new rank-$(q+2)$ transversal matroid is used through its Lorentzian basis polynomial. The theorem asserts neither real-rootedness of $p_G$ nor stability of the original complemented-row polynomial. The support identity, and not the values of nonzero determinants, determines all coefficients.
\end{remark}

\section{Exact failure at rank six and at every higher rank}\label{sec:negative}
We include the elementary obstruction calculation to make the sharpness conclusion self-contained. It is the negative construction from the Hall reports \cite{Hall,OldHall}.

Let $H_{a,b;N,M}$ have a complete $a\times b$ core, $N$ exterior-left vertices adjacent to every right-core vertex, and $M$ exterior-right vertices adjacent to every left-core vertex. Give core vertices activity $w$ and exterior vertices activity one. If $N\ge b$ and $M\ge a$, the core is a minimum cover of size $r=a+b$: match its two shores independently into the opposite exterior sets.

A size-$k$ endpoint support with $i$ selected left-core and $j$ selected right-core vertices is feasible exactly when $i+j\ge k$, subject to the cardinality bounds. Match all exterior vertices to the core first, then use the complete core for the remainder. Thus
\begin{equation}\label{eq:Hallcount}
 a_k=\sum_{\substack{0\le i\le a,\ 0\le j\le b\\i+j\ge k}}
 \binom ai\binom bj\binom N{k-i}\binom M{k-j}w^{i+j}.
\end{equation}
Put $A_i=\binom N{b-i}$, $B_i=\binom M{a-i}$ for $i=0,1,2$, and
\[
 X=aA_0B_1+bA_1B_0,\quad Y=aA_1B_2+bA_2B_1,
 \quad Z=\binom a2A_0B_2+abA_1B_1+\binom b2A_2B_0.
\]
Counting at most two omitted core vertices gives
\[
 a_r=A_0B_0w^r,\quad
 a_{r-1}=A_1B_1w^r+Xw^{r-1},\quad
 a_{r-2}=A_2B_2w^r+Yw^{r-1}+Zw^{r-2}.
\]
Hence
\begin{equation}\label{eq:neggap}
 \Delta_{r-1}^{(r)}=w^{2r-2}(\alpha w^2+\beta w+\gamma),
\end{equation}
where
\begin{align*}
 \alpha&=(r-1)A_1^2B_1^2-2rA_0A_2B_0B_2,\\
 \beta&=2(r-1)A_1B_1X-2rA_0B_0Y,\\
 \gamma&=(r-1)X^2-2rA_0B_0Z.
\end{align*}

\subsection{A concrete rank-six witness}
Take $a=b=3$ and $N=M=33$. Here $A_0=B_0=5456$, $A_1=B_1=528$, $A_2=B_2=33$, and exact simplification of \eqref{eq:neggap} gives
\[
 \Delta_5^{(6)}=404794368\,w^{10}(-w^2+26784w+522784).
\]
At $w=26804$, the parenthesis equals $-13296$. This is an exact positive-integer-activity counterexample on $72$ vertices and $207$ edges, with minimum cover and matching rank six. Combined with Corollary~\ref{cor:rankfive}, the first weighted failure rank is exactly six.

\subsection{Permanent scaling failure at every rank at least six}
Let $a,b\ge3$. Write
\[
 K=(r-1)ab,\qquad J=(r+1)ab-2r(r-1).
\]
Because $ab\ge3(r-3)$ and $r\ge6$,
\[
 J\ge r^2-4r-9>0.
\]
Choose $N=b+30$, $M=a+30$, so that
\[
 \frac{A_1^2B_1^2}{A_0A_2B_0B_2}
 =\frac{ab}{(a-1)(b-1)}\left(\frac{32}{31}\right)^2.
\]
The inequality $\alpha<0$ is equivalent to
\[
 K(2\cdot31+1)<J\cdot31^2.
\]
Indeed
\[
 15J-K=(14r+16)ab-30r(r-1)
 \ge12(r-6)(r+2)\ge0,
\]
and $63/961<1/15\le J/K$. Thus \eqref{eq:neggap} is negative for every sufficiently large $w$. An explicit sufficient integer is
\[
 w\ge1+\left\lfloor\frac{|\beta|+|\gamma|}{-\alpha}\right\rfloor,
\]
since for $w\ge1$ the bracket is at most $w(\alpha w+|\beta|+|\gamma|)<0$.
Taking $a=3$, $b=r-3$ supplies every $r\ge6$. In contrast, every graph of rank at most five remains $\ULC$ at every positive minimum-cover scale by Corollary~\ref{cor:rankfive}. This proves the sharp universal rank threshold for minimum-cover scaling.

\section{An established geometric interpretation}\label{sec:geometry}
For unit activities, $p_G$ is the \emph{perfectly matchable set polynomial} of Davis--Kohl \cite{DK}. It is distinct from the ordinary matching or rook polynomial, which counts matchings with multiplicity. Their Theorem 3.10, restating results of Ohsugi--Tsuchiya, gives the following precise Ehrhart interpretation.

Given bipartite $G=(X,Y;E)$, form $\widehat G$ by adjoining $x_0$ on the left and $y_0$ on the right, with edges from $x_0$ to every old right vertex, from $y_0$ to every old left vertex, and $x_0y_0$. For the edge polytope
\[
 Q_{\widehat G}=\operatorname{conv}\{e_x+e_y:xy\in E(\widehat G)\},
\]
one has
\[
 h^*(Q_{\widehat G};t)=p_G(t)=h^*(\operatorname{STAB}(\overline G);t).
\]
The second identity is their Lemma 3.11; $\operatorname{STAB}(\overline G)$ is the stable-set polytope of the complement.

\begin{corollary}
If $G$ has matching rank $r\le5$, or a minimum vertex cover meeting one shore in at most two vertices, these two $h^*$-polynomials are $\ULC_r$, normalized by their degree $r=\nu(G)$.
\end{corollary}
Only the unit-activity interpretation is used here. No assertion is made about arbitrary Ehrhart $h^*$-polynomials or about the $h^*$-polynomial of an unrelated preorder polytope.

\appendix

\section{A generic-line matroid model for the forced-left sector}\label{sec:line}
By \eqref{eq:union}, the polynomial in Theorem~\ref{thm:forcedleft} is
\begin{equation}\label{eq:C}
 C=Qf+a g_2+b g_1+c g_U+h.
\end{equation}
Write $f_h=s^qf(t/s)$, $G_i=s^{q-1}g_i(t/s)$, $G_U=s^{q-1}g_U(t/s)$ and $H_h=s^{q-2}h(t/s)$. For $q=1$ the last term is zero. The right-complement transversal matroid of $H+A$ gives a degree-$q$ basis specialization
\[
 F=f_h+z_1G_1+z_2G_2+z_1z_2H_h.
\]
The private $B$ variables are $s/v_b$, the exterior-left variables are $tu_l$, the two distinguished elements have variables $z_1,z_2$, and we multiply by $\prod_bv_b$.

Represent this matroid by a generic matrix: use an independent indeterminate for each allowed adjacency and private standard basis columns. A square determinant is nonzero exactly when its support admits a matching, since different matchings give distinct monomials. Denote the two distinguished columns by $v_1,v_2$. Adjoin columns
\[
 U_j=\alpha_jv_1+\beta_jv_2\quad(1\le j\le m),
\]
with generic coefficients, algebraically independent over the original matrix field. For a basis containing just one $U_j$ from this distinguished set, its determinant is $\alpha_jD_1+\beta_jD_2$ and is nonzero iff at least one of $D_1,D_2$ is nonzero. Thus its remaining-element profile is exactly $G_U$. Any pair of distinct distinguished columns has determinant a nonzero scalar multiple of the determinant using $v_1,v_2$, hence profile $H_h$. Three distinguished columns are dependent. These statements also hold when $v_1,v_2$ are dependent or zero: pair profiles then vanish, and the single profile remains the union.

The resulting matroid basis specialization is therefore
\begin{align}\label{eq:Fext}
 F_{\rm ext}={}&f_h+z_1G_1+z_2G_2+\left(\sum_jZ_j\right)G_U\notag\\
 &+\left[z_1z_2+(z_1+z_2)\sum_jZ_j+\sum_{i<j}Z_iZ_j\right]H_h.
\end{align}
It is Lorentzian. For $Q>0$, substitute
\[
 z_1=bs/Q,\qquad z_2=as/Q,\qquad Z_j=r_js/Q.
\]
The bracket becomes $s^2/Q$, and
\[
 QF_{\rm ext}=Qf_h+s(aG_2+bG_1+cG_U)+s^2H_h=s^qC(t/s).
\]
This proves the first part of Theorem~\ref{thm:forcedleft}. For $Q=0$, add two singleton right vertices with positive activities tending to zero and use coefficientwise closedness.

\begin{corollary}[Scaling only the two cover vertices]\label{cor:two-scale}
For a genuine $2+q$ cover with $q\ge1$, scaling only the two $A$ activities by a common sufficiently large factor makes every $\ULC_{q+2}$ gap strictly positive.
\end{corollary}
\begin{proof}
Every matching of size $q+2$ uses exactly one cover vertex per edge. Thus $R$ matches both $A$ vertices and $L$ matches all $B$, so $Q>0$ and $f_q>0$. Consequently all coefficients $C_0,\ldots,C_q$ are positive. Write the two activities as $w\alpha,w\beta$. With $E$ the weighted edge sum incident with $A$ at relative activities $\alpha,\beta$,
\[
 a_1(w)=wE+O(1),\qquad
 a_k(w)=w^2\alpha\beta C_{k-2}+O(w)\quad(k\ge2).
\]
The graph induced by $A$ and the opposite shore has matching rank two, so $E^2\ge4\alpha\beta C_0$. The first two leading gaps are respectively
\[
 (q+1)E^2-2(q+2)\alpha\beta C_0\ge2q\alpha\beta C_0>0,
 \qquad 2q(\alpha\beta C_0)^2>0,
\]
at powers $w^2,w^4$. For $k=i+2$, $1\le i<q$, the leading coefficient at $w^4$ is at least
\[
 \frac{2(q-i+1)}{i}(\alpha\beta)^2C_{i-1}C_{i+1}>0,
\]
by substituting the $\ULC_q$ inequality for $C_i^2$ into \eqref{eq:gap}. Positive leading coefficients prove the claim.
\end{proof}

\section{The sharp all-core inequality at rank five}\label{sec:forcedcore}
The full lift gives a short general boundary theorem before the independent rank-five certificate below.
\begin{theorem}[Sharp all-core constant for a two-vertex shore]\label{thm:generalcore}
For a $2+q$ cover with $q\ge2$, factor the product of all core activities from the all-core sector and write it as $t^q(c_0+c_1t+c_2t^2)$. Then
\[
 c_1^2\ge\frac{2q}{q-1}c_0c_2,
\]
and the constant is optimal for every $q$.
\end{theorem}
\begin{proof}
In the lifted matroid basis polynomial, set all private $B$ variables to zero and differentiate once in each augmented core variable $X,Y$. These operations preserve the Lorentzian property, giving a degree-$q$ polynomial. Under the remaining specializations its three plane-count sectors give, up to the fixed positive scaling factor,
\[
 c_0s^2t^{q-2}+c_1st^{q-1}+c_2t^q.
\]
The bivariate $\ULC_q$ inequality at index $q-1$ is the desired bound. Zero activities and $Q=0$ follow by continuity as before.

For the complete $2\times q$ core with $N\ge q$ unit common-left vertices and $M\ge2$ unit common-right vertices,
\[
 c_0=\binom N{q-2},\qquad c_1=M\binom N{q-1},\qquad
 c_2=\binom M2\binom Nq.
\]
Their ratio is
\[
 \frac{c_1^2}{c_0c_2}
 =\frac{2q}{q-1}\frac{M}{M-1}\frac{N-q+2}{N-q+1}
 \longrightarrow\frac{2q}{q-1}.
\]
\end{proof}

The following argument independently proves the $q=3$ case using only a finite quartic coefficient test and explicit sums of squares. It does not rely on the full lift.

Now fix a $2+3$ cover. Strip off the product of the five core activities from the sector in which all five core vertices are selected. Its polynomial is
\[
 t^3(c_0+c_1t+c_2t^2).
\]
\begin{theorem}[Sharp constant three]\label{thm:three}
For arbitrary nonnegative exterior activities and every $2+3$ core,
\[
 c_1^2\ge3c_0c_2.
\]
The constant $3$ cannot be increased uniformly.
\end{theorem}

Let $T$ be the weight of exterior-left triples matching all of $B$. Let $C$ be the weight of single exterior-left vertices which, together with the two core rows, match $B$. Let $P_i$ be the weight of exterior-left pairs completable with core row $S_i$, and let $P_U$ use the union row. Set
\[
 O=P_1+P_2-P_U.
\]
The union interpretation gives $0\le O\le\min(P_1,P_2)$. The central left-side inequality is
\begin{equation}\label{eq:leftkey}
 4P_1P_2-2O^2\ge3CT.
\end{equation}

\subsection{A finite exact coefficient certificate}
Encode a nonempty neighborhood in $B$ by its bit mask $1,\ldots,7$. For three nonempty masks, let $M(r,s,t)$ be one exactly when they admit a matching to all three columns. Hall's condition gives the elementary rule: their bitwise union is $7$, and no singleton mask $1,2,4$ appears twice. For general tuples including an empty mask, use ordinary Boolean matchability.

Index left exterior vertices by $i$, with masks $t_i$ and weights $z_i$. Put
\begin{align*}
 c_i&=M(S_1,S_2,t_i),&p_{ij}&=M(S_1,t_i,t_j),&q_{ij}&=M(S_2,t_i,t_j),\\
 o_{ij}&=p_{ij}q_{ij},&\tau_{ijk}&=M(t_i,t_j,t_k).
\end{align*}
Then $C=\sum c_iz_i$, $P_1=\sum_{i<j}p_{ij}z_iz_j$, and similarly for $P_2,O,T$. In the quartic
\[
 I=4P_1P_2-2O^2-3CT
\]
every exponent is at most two. Its coefficient of $z_i^2z_j^2$ is $2o_{ij}\ge0$. Its coefficient of $z_i^2z_jz_k$ for distinct indices is
\begin{equation}\label{eq:repeatcoef}
 4(p_{ij}q_{ik}+p_{ik}q_{ij})-4o_{ij}o_{ik}-3c_i\tau_{ijk}.
\end{equation}
Its coefficient of $z_1z_2z_3z_4$ is
\begin{equation}\label{eq:squarefreecoef}
 4\sum_{ij\mid kl}(p_{ij}q_{kl}+p_{kl}q_{ij})
 -4\sum_{ij\mid kl}o_{ij}o_{kl}
 -3\sum_i c_i\tau_{[4]\setminus\{i\}},
\end{equation}
where each of the three unordered pair partitions is used once.

If the core has matching rank at most one, $C=0$ and \eqref{eq:leftkey} follows from $O^2\le P_1P_2$. Otherwise row exchange and column permutation give exactly the eight representatives in Table~\ref{tab:coeff}. The table records the complete range of possible coefficients in \eqref{eq:repeatcoef} and \eqref{eq:squarefreecoef}. It is obtained by the stated three-row Boolean rule, over $7^3=343$ ordered triples and $\binom{10}{4}=210$ neighborhood multisets per core. The included exact verifier implements an independent permutation matcher and checks all $4424$ substitutions, as well as all $46$ labeled matching-rank-two cores.

\begin{table}[ht]
\centering
\begin{tabular}{ccc}
\toprule
Core masks & Repeated-exponent coefficients & Squarefree coefficients\\
\midrule
$(1,2)$ & $0,1,4$ & $0,2,3$\\
$(1,3)$ & $0,1,4$ & $0,2,3,6,8,12$\\
$(1,6)$ & $0,1,4$ & $-4,0,2,3$\\
$(1,7)$ & $0,1,4$ & $-4,0,2,3,8,12$\\
$(3,3)$ & $0,1,4$ & $0,2,3,6,8,12$\\
$(3,5)$ & $0,1,4$ & $0,2,3,8,12$\\
$(3,7)$ & $0,1,4$ & $0,2,3,8,12$\\
$(7,7)$ & $0,1,4$ & $0,2,3,8,12$\\
\bottomrule
\end{tabular}
\caption{The complete finite coefficient test. In each exceptional row the only four negative neighborhood multisets are listed in \eqref{eq:bad}.}\label{tab:coeff}
\end{table}

Only cores $(1,6)$ and $(1,7)$ have negative coefficients. In both cases the negative neighborhood multisets are exactly
\begin{equation}\label{eq:bad}
 (2,3,4,5),\quad(2,3,5,5),\quad(3,3,4,5),\quad(3,3,5,5),
\end{equation}
and each coefficient is $-4$. Thus every term involving another neighborhood class is nonnegative. It remains to treat the four classes $2,3,4,5$ in these two cores.

\subsection{The four-class sum-of-squares certificate}
For core $(1,6)$ let $a,b,c,d$ now denote the activity sums in classes $2,3,4,5$, and let $B_2,D_2$ be the second elementary moments of classes $3,5$. Put $u=a+b$, $v=c+d$. Exact support counting yields
\begin{align*}
 C&=u+v,& P_1&=uv,\\
 P_2&=(a+c)(b+d)+bd+B_2+D_2,& O&=ad+bc+bd,\\
 T&=v(ab+B_2)+u(cd+D_2).
\end{align*}
Since $0\le B_2\le b^2/2$ and $0\le D_2\le d^2/2$, and $I$ is affine in both moments, it suffices to check the four rectangle corners. Write $I_{00},I_{01},I_{10},I_{11}$ for the lower/lower, lower/upper, upper/lower and upper/upper choices. The following exact identities have only nonnegative summands:
\begin{align*}
 I_{00}={}&2(ad-bc)^2+uv(ab+cd)+(ac+bd)(ad+bc)+2b^2d^2,\\
 I_{01}={}&\tfrac12d^2(a-b)^2+\tfrac12ad(b-d)^2
          +ac(b-d)^2+bc(a-d)^2+R_{01},\\
 R_{01}={}&a^2(bd+cd)
 +a\left(\tfrac12b^2d+bc^2+c^2d+\tfrac12cd^2\right)\\
 &+b^2(2c^2+cd)+b\left(c^2d+\tfrac12cd^2+\tfrac12d^3\right),\\
 I_{11}={}&\tfrac12bd(u-v)^2+\tfrac12(ad+bc+3ac)(b-d)^2\\
 &+(b+d)\left(ac(a+c)+\tfrac12(a^2d+bc^2)\right).
\end{align*}
The expression for $I_{10}$ is obtained from $I_{01}$ by interchanging $(a,b)$ with $(c,d)$. For core $(1,7)$, the restricted $P_2,O$ both increase by $ac$ while $C,T,P_1$ are unchanged. Since $P_1-O=ac$, its quartic is exactly the preceding $I+2a^2c^2$. This proves \eqref{eq:leftkey} in every case. The included symbolic checker verifies all five identities exactly.

\subsection{Passing to the right exterior and proving sharpness}
Return to the right-side notation $a,b,c,d$ of Section 3. Selecting all core vertices gives
\[
 c_0=C,\qquad c_1=aP_2+bP_1+cP_U,\qquad c_2=TQ.
\]
Let $A=a+c$, $B=b+c$, $p=P_1$, $q=P_2$. Then
\[
 c_1=qA+pB-Oc,\qquad Q\le AB-c^2/2.
\]
For $pq>0$, the exact quadratic identity
\begin{align*}
 &(qA+pB-Oc)^2-(4pq-2O^2)(AB-c^2/2)\\
 &\quad=\left(1-\frac{O^2}{2pq}\right)(qA-pB)^2
 +\frac{[O(qA+pB)-2pqc]^2}{2pq}\ge0
\end{align*}
uses only $O^2\le pq$. Together with \eqref{eq:leftkey} it implies $c_1^2\ge3CTQ=3c_0c_2$. If $pq=0$, the same assertion follows by continuity, or directly from $O=0$ and $CT=0$.

For sharpness, take the complete $2\times3$ core, $N$ common left-exterior vertices and $M$ common right-exterior vertices, all with unit exterior activities. Then
\[
 c_0=N,\quad c_1=\binom N2M,\quad c_2=\binom N3\binom M2,
 \qquad
 \frac{c_1^2}{c_0c_2}=3\frac{N-1}{N-2}\frac{M}{M-1}\longrightarrow3.
\]
This completes Theorem~\ref{thm:three}.

\section{Lorentzian transfer across a star core}\label{sec:stars}
\begin{theorem}[Star-core transfer]\label{thm:stars}
For a displayed $2+q$ cover with matching-rank-at-most-one core, both root orientations admit direct Lorentzian transfer proofs of $\ULC_{q+2}$.
\end{theorem}
This is now a special case of Theorem~\ref{thm:main}. We retain the independent proofs because the transfer operators and their interlacing certificates are useful beyond the full-lift construction.
\subsection{A support-exact vertex-sum identity}
Suppose $G_1,G_2$ intersect in just one vertex $v$ of activity $u$, and write
\[
 p_1=A(t)+utB(t),\qquad p_2=C(t)+utD(t).
\]
Then
\begin{equation}\label{eq:vertices}
 p_{G_1\cup G_2}=AC+ut(BC+AD).
\end{equation}
If $v$ is absent, each component has balanced selected shores. If $v$ is selected, exactly one component uses it, and that component is determined by its selected-vertex imbalance away from $v$: its opposite shore has one extra selected vertex. Thus the two root-using terms cannot count the same endpoint support, even if that support has several perfect matchings.

\subsection{Stars centered on the two-vertex shore}
Suppose every core edge is incident with $v=A_1$. Let $G_1$ have shores $L\cup\{v\}$ and $B$, and let $G_2$ have shores $\{v,A_2\}$ and $R$. Their intersection is just $v$. The right-complement basis construction for $G_1$ gives a Lorentzian polynomial
\[
 F(s,t,z)=A_h(s,t)+zB_h(s,t),
 \quad A_h=s^qA(t/s),\quad B_h=s^{q-1}B(t/s).
\]
Here $z$ is the root element variable; setting $z=ut$ gives the degree-$q$ homogenization of $p_1$.

For $G_2$, write $C=1+dt$, $D=e+ft$. If $A_2$ has activity $y$, and the exterior class sums are $a,b,c$ with common second moment $\delta$, then
\[
 d=y(b+c),\quad e=a+c,\quad f=y(ab+ac+bc+\delta),
 \qquad de-f=y(c^2-\delta)\ge0.
\]
The homogeneous quadratic
\[
 K(s,t,z)=s^2+dst+esz+ftz
\]
is Lorentzian. Its Hessian is
\[
 \begin{pmatrix}2&d&e\\d&0&f\\e&f&0\end{pmatrix},
 \qquad \det=2f(de-f)\ge0.
\]
For $f>0$, its $t,z$ principal submatrix has one positive and one negative eigenvalue; the determinant and eigenvalue interlacing exclude two positive eigenvalues of the full Hessian. For $f=0$, $K=s(s+dt+ez)$ is a product of nonnegative linear forms.

The product $FK$ is Lorentzian. Retain only its terms of degree at most one in $z$, then set $z=ut$. The output is
\[
 A_h(s^2+dst)+ut\big[B_h(s^2+dst)+A_h(es+ft)\big],
\]
which is the degree-$(q+2)$ homogenization of \eqref{eq:vertices}. Partial truncation preserves Lorentzian polynomials: polarize $s,t$, take the full multiaffine part (only powers $z^2$ are then removed), and identify the polarized variables again. Thus the output is Lorentzian and $p_G$ is $\ULC_{q+2}$.

\subsection{Stars centered on the opposite shore}
Suppose every core edge is incident with $v=B_1$. Let $G_1=G[L,B]$ and let $G_2$ have shores $A$ and $R\cup\{v\}$. For $G_1$, the private-root variable now gives
\[
 F(s,t,z)=zA_h+tB_h,
 \qquad A_h=s^{q-1}A(t/s),\quad B_h=s^{q-1}B(t/s).
\]
The other private variables are $s/v_b$, the left variables are $tu_l$, and the basis polynomial is multiplied by $\prod_{b\ne v}v_b$. Hence $F$ is Lorentzian and $uF(s,t,s/u)=s^qp_1(t/s)$.

For $G_2$ put $C_h=s^2+c_1st+c_2t^2$ and $D_h=es+ft$. Give its two left vertices fixed activities $\alpha,\beta$; these are constants, distinct from the private-root indeterminate $z$. With exterior class sums $a,b,c$ and second moment $\delta$,
\[
 c_1=\alpha(a+c)+\beta(b+c),\qquad c_2=\alpha\beta(ab+ac+bc+\delta),
\]
so
\[
 c_1^2-4c_2=[\alpha(a+c)-\beta(b+c)]^2+4\alpha\beta(c^2-\delta)\ge0.
\]
The rooted interlacing inequality is
\begin{equation}\label{eq:rootinterlace}
 c_1ef-f^2-c_2e^2\ge0.
\end{equation}
If $v$ is adjacent only to the first left vertex, $e=\alpha$, $f=\alpha\beta(b+c)$ and the expression is $\alpha^3\beta(c^2-\delta)$. The other singleton case gives $\alpha\beta^3(c^2-\delta)$. If $v$ is adjacent to both, $e=\alpha+\beta$, $f=\alpha\beta(a+b+c)$ and it is exactly
\[
 \alpha\beta\left[
 \left(a\alpha-b\beta+\tfrac c2(\alpha-\beta)\right)^2
 +\left(\tfrac{3c^2}{4}-\delta\right)(\alpha+\beta)^2\right]\ge0.
\]
An isolated root has $e=f=0$.

Factor $C_h=(s+\lambda_1t)(s+\lambda_2t)$, with $0\le\lambda_1\le\lambda_2$. If $e>0$, \eqref{eq:rootinterlace} gives $\lambda_1\le f/e\le\lambda_2$. Therefore
\[
 h(s,t)=\frac{stD_h}{C_h}
\]
is a nonnegative linear combination of $st/(s+\lambda_i t)$ with positive total coefficient. Each such fraction maps $\HH^2$ into $\HH$, since its reciprocal is $1/t+\lambda_i/s$, of strictly negative imaginary part. Repeated $\lambda_i$ give the same conclusion directly.

Consider the homogeneous degree-two operator
\[
 \mathcal T(P)=C_hP+stD_h\partial_zP.
\]
It preserves real stability. At upper-half-plane values, a stable nonzero $P$ has $\operatorname{Im}(\partial_zP/P)\le0$, by its univariate factorization in $z$. A zero of $\mathcal T(P)$ with $D_h\ne0$ would require this quotient to equal $-1/h$, which has strictly positive imaginary part. If $D_h=0$, stable multiplication suffices.

The operator also preserves nonnegative coefficients. Restrict it to any finite coordinatewise degree box containing $F$; its output lies in a larger degree box. Br\"and\'en--Huh's Theorem 3.4 applies, so $\mathcal T(F)$ is Lorentzian. Finally
\[
 u\mathcal T(F)|_{z=s/u}=sA_hC_h+utB_hC_h+ustA_hD_h
\]
is precisely the degree-$(q+2)$ homogenization of \eqref{eq:vertices}. This proves the column-star case.

Every nonempty bipartite graph of matching rank one has all its edges incident with a common vertex. The two orientations therefore exhaust such cores. An empty core gives a product of $\ULC_q$ and $\ULC_2$ polynomials, which is $\ULC_{q+2}$ by Lorentzian product closure. Theorem~\ref{thm:stars} follows.

\section{Exact equality in the first Newton gap}\label{sec:firstgap}
The following gives a useful structural interpretation of non-strictness. Its ingredient is the classical Motzkin--Straus bound \cite{MS}; we include the short argument in the form needed here.

\begin{theorem}
For a positive weighted bipartite graph of matching rank $r\ge2$,
\[
 (r-1)a_1^2-2ra_2=0
\]
if and only if every nontrivial connected component is a star and their weighted edge totals are all equal. In that case $p_G(t)=(1+Et)^r$ for their common edge total $E$; otherwise the first gap is strict.
\end{theorem}
\begin{proof}
Form the compatibility graph $K$ whose vertices are the edges of $G$, with two vertices adjacent when the corresponding edges are disjoint. Give an edge $xy$ weight $w_{xy}=u_xv_y$. Put $S=\sum w_e=a_1$ and
\[
 m_2=\sum_{\{e,f\}\in E(K)}w_ew_f.
\]
The clique number of $K$ is $r$, and $m_2\ge a_2$ because $m_2$ counts two-edge matchings rather than endpoint supports.

For completeness, the Motzkin--Straus bound is
\[
 m_2\le\frac{r-1}{2r}S^2.
\]
To see it, merge the weights of two nonadjacent vertices onto whichever has the larger weighted neighbor sum. This preserves $S$ and does not decrease $m_2$. Continue until the positive support is a clique of size at most $r$, where the bound follows from Cauchy--Schwarz.

Equality in the asserted first gap forces equality in both preceding bounds. Since the original compatibility weights are all positive, they give an interior maximum of $m_2$ on the simplex of fixed sum. If $K$ were not complete multipartite, it would contain an induced edge $ik$ and a vertex $j$ nonadjacent to both. The sum-preserving direction $e_i-2e_j+e_k$ has strictly positive second derivative for $m_2$, contradicting a local maximum. Hence $K$ is complete multipartite. Its complement, the line graph of $G$, is a disjoint union of cliques. In a bipartite graph a nontrivial connected component whose edges are pairwise intersecting is a star. There are exactly $r$ such components. The compatibility bound is then equality precisely when the total weights in the $r$ parts are equal. Conversely, for stars support and matching counts coincide and the polynomial is the stated product.
\end{proof}

\section{Quantitative eventual strictness under common cover scaling}\label{sec:scaling}
Fix a genuine $2+3$ cover whose core has matching rank two. Write every cover activity as $w$ times its fixed positive relative activity. Let $e_1,e_2$ be the first two support coefficients of the induced core at those relative activities, and let $\zeta$ be the product of all five relative activities. For the all-core coefficients of Theorem~\ref{thm:three}, set $\gamma_i=\zeta c_i$.

Every size-five matching uses exactly one cover vertex per edge; a core edge would consume two cover vertices and preclude five edges. Thus the exterior-left graph matches all three $B$ vertices and the exterior-right graph matches both $A$ vertices, giving $c_2>0$. A core two-edge matching leaves one $B$ vertex unmatched; that vertex has an exterior-left neighbor, so adding it gives $c_0>0$. Also $e_1,e_2>0$ and $e_1^2\ge4e_2$ by Lemma~\ref{lem:small} for the core.

Counting selected cover vertices gives
\begin{align*}
 a_0&=1,&a_1&=e_1w^2+O(w),&a_2&=e_2w^4+O(w^3),\\
 a_3&=\gamma_0w^5+O(w^4),&a_4&=\gamma_1w^5+O(w^4),&a_5&=\gamma_2w^5.
\end{align*}
Consequently the leading gap limits are
\begin{align*}
 \Delta_1^{(5)}/w^4&\longrightarrow4e_1^2-10e_2\ge6e_2>0,\\
 \Delta_2^{(5)}/w^8&\longrightarrow6e_2^2>0,\\
 \Delta_3^{(5)}/w^{10}&\longrightarrow6\gamma_0^2>0,\\
 \Delta_4^{(5)}/w^{10}&\longrightarrow4\gamma_1^2-10\gamma_0\gamma_2
 \ge2\gamma_0\gamma_2>0.
\end{align*}
The last inequality uses the sharp constant three from Theorem~\ref{thm:three}. All four gaps are therefore strictly positive eventually. This is stronger than the already proved non-strict $\ULC_5$ property at every $w>0$.

For either this scaling or Corollary~\ref{cor:two-scale}, write each gap as
\[
 \Delta_k(w)=L_kw^{d_k}+\sum_{j<d_k}b_{kj}w^j,\qquad L_k>0.
\]
A directly computable sufficient threshold is
\begin{equation}\label{eq:positivethreshold}
 W=1+\max_k\frac{\sum_{j<d_k}\max(-b_{kj},0)}{L_k}.
\end{equation}
For $w\ge W$, every negative lower term is bounded in magnitude by its coefficient times $w^{d_k-1}$, so every gap is positive. This threshold depends on the graph and its fixed field; no uniform bound is asserted.

For a rank-five minimum cover of split $0+5$ or $1+4$, Lemma~\ref{lem:singleton} already gives $\ULC_5$ at every scale. For a $2+3$ cover, Theorem~\ref{thm:main} does likewise; the calculation above adds strictness when its core has matching rank two. Thus all choices of minimum cover are covered, with the relative activities fixed before scaling.

\section{Exact verification and scope of the conclusions}\label{sec:verification}
The proofs above are universal arguments. Their exact checks are independently implementable and do not use numerical root searches or numerical optimization as evidence for a theorem.
\begin{itemize}
\item The full transversal lift is checked by Boolean basis enumeration on $180$ weighted graphs with $q=0,1,2,3,4,5$. It tests $140450$ candidate basis subsets and compares the complete left-subset marginal, before diagonalizing the left variables.
\item An independent finite-field representation check covers all $64$ labeled $2\times3$ cores and $73441$ basis determinants, comparing the full weighted support polynomial exactly. A separate generic-line check covers $192$ graphs and $23553$ determinants, including $54$ rank-deficient cores.
\item The union formula is checked on $1352$ weighted graphs with $q=1,2,3,4$, including overlapping, nested and empty core neighborhoods.
\item The rank-five quartic certificate checks all $4424$ stated coefficient cases, all $46$ labeled rank-two cores, the four moment-corner identities, the exceptional correction, and the right-side quadratic identity.
\item The star-core and scaling checks cover $820$ weighted star-core graphs for $q=1,\ldots,5$ and $230$ rank-five common-scaling instances. All displayed interlacing sums of squares, grading identities and computed strictness thresholds are verified exactly.
\item The negative-family verifier checks complete endpoint support counts for seven core shapes and the universal exterior population choice for all $3\le a,b\le20$. The rank-six witness and its exact negative gap are included.
\end{itemize}
The package contains the Python sources, fixed seeds and JSON outputs. Only SymPy is needed beyond the Python standard library. Run \texttt{python3 checks/run\_all.py}; the JSON outputs are regenerated in \texttt{data/}. The LaTeX sources compile with ordinary pdfLaTeX; \texttt{build\_local.sh} also handles the local TeX configuration used to generate this PDF.

R\"ohrle--Ulirsch's bimatroid theorem \cite{RU} normalizes by the smaller ambient shore size. The new lift replaces that order by $q+2$, and hence by matching rank when the cover is minimum. The polynomial is the perfectly matchable set polynomial of \cite{DK}, not a matching count with multiplicities. A bounded primary-literature comparison located no prior universal $2+q$-cover or weighted rank-five theorem; global priority remains unclaimed.

Higher-rank unrestricted unit-weight and one-shore-only questions are outside this report. The proof uses Lorentzian matroid bases; it does not assume that transversal matroids have the half-plane property. The cover obstruction in Corollary~\ref{cor:coverobstruction} is sharp.

\begin{thebibliography}{10}
\small
\setlength{\itemsep}{2pt}
\bibitem{BH}
P. Br\"and\'en and J. Huh, \emph{Lorentzian polynomials}, Annals of Mathematics \textbf{192} (2020), 821--891.
\href{https://arxiv.org/html/1902.03719v3}{arXiv:1902.03719v3}. Theorem 2.10, Example 2.26, Corollary 2.32, Proposition 3.1, Theorem 3.4, Corollary 3.5 and Theorem 3.10.

\bibitem{RU}
F. R\"ohrle and M. Ulirsch, \emph{Logarithmic concavity of bimatroids}, Annals of Combinatorics (2025).
\href{https://arxiv.org/html/2402.15317v3}{arXiv:2402.15317v3}. Theorem A and Example 2.4.

\bibitem{DK}
R. Davis and F. Kohl, \emph{Perfectly Matchable Set Polynomials and $h^*$-polynomials for Stable Set Polytopes of Complements of Graphs}, Discrete Mathematics \textbf{347}(7) (2024), 114018.
\href{https://arxiv.org/html/2207.14759}{arXiv:2207.14759}. Theorem 3.10 and Lemma 3.11 state the geometric identities used here.

\bibitem{MS}
T. S. Motzkin and E. G. Straus, \emph{Maxima for Graphs and a New Proof of a Theorem of Tur\'an}, Canadian Journal of Mathematics \textbf{17} (1965), 533--540.
\href{https://doi.org/10.4153/CJM-1965-053-6}{doi:10.4153/CJM-1965-053-6}.

\bibitem{Rank}
\emph{Rank-Three Ultra-Log-Concavity and Exact Kernels for Matching Supports}, ProveIt research manuscript prepared for Vladimir Reshetnikov with ChatGPT, 30 September 2026.
\href{https://github.com/VladimirReshetnikov/ProveIt/blob/1b3960d8acdabd983147b7c609bd9dfb2d67a0a0/docs/incoming/ProveIt_Matching_Rank_Research.zip}{Pinned incoming archive}.

\bibitem{Leaf}
\emph{Leaf Compression and Weighted Rank Three for Matching Support Polynomials}, companion research note prepared for Vladimir Reshetnikov with OpenAI, 1 October 2026. The compression proof is repeated in Lemma~\ref{lem:singleton}.

\bibitem{Four}
\emph{Stability of Two by Two Covers and Weighted Rank Four}, companion research note prepared for Vladimir Reshetnikov with OpenAI, 1 October 2026. The present rank-five ULC proof is self-contained and does not import its stability constructions.

\bibitem{Hall}
\emph{Sharp Full Hall Classification}, companion research report prepared for Vladimir Reshetnikov with OpenAI, 1 October 2026. Its general negative-shape construction is reproduced in Section~\ref{sec:negative}.

\bibitem{OldHall}
\emph{Hall Bottlenecks and the Limits of Rank Normalization}, ProveIt research manuscript prepared for Vladimir Reshetnikov with ChatGPT, 30 September 2026.
\href{https://github.com/VladimirReshetnikov/ProveIt/blob/a866ff9a2/docs/incoming/ProveIt_Hall_Bottlenecks.zip}{Pinned incoming archive}. The rank-six witness is restated and verified here.
\end{thebibliography}

\end{document}
